\subsubsection{Uniform ramification breaks from the different}

\begin{lemma}[Proved: a degree-independent tube bound]
\label{D:different-conductor-bound}
Let $A$ be a complete discrete valuation ring of positive residue
characteristic, with fraction field $K$
and uniformizer $\pi$, and let $B$ be its integral closure in a finite
separable field extension $L/K$. Normalize $v$ by $v(\pi)=1$.
If $\pi^M\Omega_{B/A}=0$ for an integer $M\geq0$, then the
nonlogarithmic Abbes--Saito conductor satisfies
\[
                         c(L/K)\leq2M.                 \tag{D.UB.1}
\]
In particular, writing $e=e(L/K)$ and
$\delta=v_L(\mathfrak D_{B/A})$, we have
\[
                 c(L/K)\leq2\left\lceil\delta/e\right\rceil .
                                                        \tag{D.UB.2}
\]
No separability of the residue extension is assumed.
\end{lemma}

\begin{proof}
We use the tube definition in \cite[sections 3.1 and 6,
Proposition 6.2(i), Definition 6.3 and Proposition 6.4,
preprint pp.5--6 and 15--16]{AbbesSaito2002D}.
Lemma 7.1 of that source, on preprint p.19, gives a presentation
\[
                 B=A[X_1,\ldots,X_r]/(f_1,\ldots,f_r).
\]
Let $J=(\partial f_i/\partial X_j)$. Its transpose presents
$\Omega_{B/A}$. The annihilation assumption therefore gives
a matrix $H$ over $B$ with $JH=\pi^M I_r$.
Indeed, transposing the initial relation gives a scaled left
inverse of $J$, and over the field $L$ this is also a scaled
right inverse. Integral polynomial lifts of $H$ satisfy
\[
                  JH-\pi^M I_r=\sum_i f_i H_i             \tag{D.UB.3}
\]
for polynomial matrices $H_i$ over $A$.

Fix a rational $a>2M$, and work in a complete algebraically
closed valued extension of $K$. On the tube
\[
 X^a=\{x:\ v(x_j)\geq0,\ v(f_i(x))\geq a\text{ for all }i,j\},
\]
(D.UB.3) gives $J(x)H(x)=\pi^M(I_r+E(x))$, with
$v(E(x))\geq a-M>0$. Thus $J(x)$ is invertible and every
entry of its inverse has valuation at least $-M$.
Here valuations of vectors and matrices mean the minimum
of the valuations of their entries.

Apply Newton iteration
$x_{n+1}=x_n-J(x_n)^{-1}f(x_n)$, starting at $x_0=x$.
If $v(f(x_n))\geq a_n$, the correction has valuation
at least $a_n-M$, and integral Taylor expansion gives
$v(f(x_{n+1}))\geq2(a_n-M)$.
With $a_0=a$ and $a_{n+1}=2(a_n-M)$, we have
$a_n=2M+2^n(a-2M)\to+\infty$. Every iterate remains integral
and in $X^a$. Completeness gives a zero $\alpha$ of all the
$f_i$ such that
\[
                          v(x-\alpha)\geq a-M>M.          \tag{D.UB.4}
\]
Since $B[1/\pi]=L$, this zero is one of the finitely many
embeddings into the algebraic closure. We do not assert
convergence inside the noncomplete separable closure.

Distinct zeros $\alpha,\beta$ satisfy
$v(\alpha-\beta)\leq M$. Otherwise put $h=\beta-\alpha$,
$b=v(h)>M$. Taylor expansion gives
$h=-J(\alpha)^{-1}Q(h)$ with $v(Q(h))\geq2b$, which would
force $b\geq2b-M>b$, a contradiction.

Over a finite splitting extension containing elements of
the required rational valuations, intersect $X^a$ with
the finitely many closed polydiscs
\[
                    v(x_j-\alpha_j)\geq a-M\quad(1\leq j\leq r).
\]
These rational subdomains are pairwise disjoint, and
(D.UB.4) says that they cover all geometric points.
They form a finite admissible partition, each member
open and closed in the tube. No component contains two
zeros. Proposition 6.2(i) of the cited source says that
every geometric component contains a zero. Hence the
embeddings map bijectively to the components. Ramification
is bounded by every $a>2M$, which proves (D.UB.1), including
$M=0$.

For a finite flat complete intersection the different is
the zeroth Fitting ideal of $\Omega_{B/A}$ and annihilates it
\cite[section 49.7 and Lemma 49.12.3, Tags 0BVV and 0BWB]{StacksDDifferent}.
If $M=\lceil\delta/e\rceil$, then $v_L(\pi^M)=eM\geq\delta$,
so $\pi^M$ belongs to the different. This proves (D.UB.2).
\end{proof}

\begin{proposition}[Proved: a uniform different bound on a fixed surface resolution]
\label{D:surface-uniform-different}
Let $(S,s)$ be a normal surface germ of finite type over an
algebraically closed field. Let $Y\to S$ be a projective
resolution which is an isomorphism away from $s$, with
exceptional integral curves $E_1,\ldots,E_r$.
Choose positive integers $a_i$ such that
\[
        A=\sum_i a_iE_i,\qquad c_i=-A\cdot E_i>0.           \tag{D.UB.5}
\]
For any connected finite Galois cover of $S-\{s\}$,
let $e_i,\delta_i$ be the ramification index and different
exponent at any prolongation of $E_i$. Then
\[
           \sum_i c_i\,\delta_i/e_i
                         \leq K_Y\cdot A+\sum_i a_i.       \tag{D.UB.6}
\]
The same statement holds for covers of the punctured
strict henselization, with the same intersection numbers.
\end{proposition}

\begin{proof}
The exceptional intersection matrix is negative definite
with nonnegative off-diagonal entries
\cite[section 10.1, Theorem 10.1]{Kollar2013D}.
Put $M=-(E_i\cdot E_j)$ and solve $Ma=(1,\ldots,1)^t$.
To see that $a$ is nonnegative, write $a=a_+-a_-$ with
disjointly supported nonnegative vectors. If $a_-\ne0$, then
$a_-^tMa=a_-^tMa_+-a_-^tMa_-\leq-a_-^tMa_-<0$,
using the off-diagonal signs and positive definiteness.
This contradicts $a_-^tMa=a_-^t(1,\ldots,1)^t>0$.
Clear denominators.
Every coordinate is positive: if one were zero, its row
equation would have nonpositive left side.
This constructs (D.UB.5).

Normalize $S$ in the cover, normalize $Y$ in the same field,
and resolve the latter surface, preserving the smooth open
away from the exceptional locus. We obtain a smooth surface
$Z$, proper over the normal source $S'$, and a proper
generically finite morphism $f:Z\to Y$ of degree $n$
\cite[Theorem 54.14.5, Tag 0BGP]{StacksDSurfaceResolution}.
Contract only smooth rational $(-1)$-curves which map to
points of $Y$, using
\cite[Lemma 54.16.8, Tag 0C2M]{StacksDSurfaceResolution}.
The morphism to $Y$ descends through each contraction:
its graph image over the contracted smooth surface is
proper and quasi-finite, since its fibre over the contracted
point is one point. The projection is finite birational onto
a normal surface, hence is an isomorphism.
Each contraction decreases the number of exceptional
components, so this process terminates.

Let $F_{i\lambda}$ dominate $E_i$, and let $Q_j$ denote the
remaining exceptional curves. All are exceptional over
the finite set above $s$ in $S'$, so each has negative square.
Adjunction gives
\[
        K_Z\cdot Q_j=2p_a(Q_j)-2-Q_j^2\geq0,\qquad
                          K_Z\cdot F_{i\lambda}\geq-1.    \tag{D.UB.7}
\]
For $Q_j$ the only possible negative case would be a
smooth rational $(-1)$-curve, which has been removed.
For $F_{i\lambda}$ use only $p_a\geq0$ and $F_{i\lambda}^2\leq-1$.

The divisor $f^*A$ is effective, has coefficient
$a_i e_{i\lambda}$ along $F_{i\lambda}$, and has nonnegative
coefficients along all the $Q_j$. Consequently
\[
       K_Z\cdot f^*A
          \geq-\sum_i a_i\sum_\lambda e_{i\lambda}
          \geq-n\sum_i a_i,                              \tag{D.UB.8}
\]
because $\sum_\lambda e_{i\lambda}f_{i\lambda}=n$ with
$f_{i\lambda}$ the full residue degree.

The determinant of the differential defines an effective
divisor $R_f$ with $K_Z=f^*K_Y+R_f$. The morphism is
\'etale outside the exceptional locus. At $F_{i\lambda}$
its coefficient is the different exponent
$\delta_{i\lambda}$, also for inseparable residue extensions,
by \cite[Lemma 49.12.6, Tag 0BWJ]{StacksDDifferent}.
Here $f$ is quasi-finite on a neighborhood of the generic
point of each $F_{i\lambda}$, which is where that lemma is
applied; the global effective divisor only uses the
nonzero determinant of the differential.
Every point-mapped curve has intersection zero with $f^*A$.
Projection therefore gives the exact identity
\[
    K_Z\cdot f^*A
       =nK_Y\cdot A-
           \sum_i c_i\sum_\lambda\delta_{i\lambda}f_{i\lambda}.
                                                               \tag{D.UB.9}
\]
Galois conjugacy of the prolongation DVRs makes
$e_{i\lambda},\delta_{i\lambda}$ independent of $\lambda$,
so the inner sum is $n\delta_i/e_i$.
Combining (D.UB.8) and (D.UB.9) proves the assertion.
No equivariant resolution is required: the conjugacy is
intrinsic to the valuations, all of which survive on $Z$.

A finite cover and its finite Galois action over the punctured
strict henselization descend to an \'etale neighborhood.
The exceptional curves and their intersection numbers
are unchanged, since the residue field is algebraically
closed. Applying the same argument there proves the
strict-local assertion.
\end{proof}

\begin{proposition}[Proved: the original fixed-boundary correction]
\label{D:fixed-boundary-uniform-different}
In Proposition~\ref{D:surface-uniform-different}, let $T\subset S$
be any proper closed subset containing $s$, and let $D$ be
the strict transform on $Y$ of its reduced divisorial part.
For every connected finite Galois cover of $S-T$ which is
tame at the original codimension-one points of $T$, we have
\[
         \sum_i c_i\,\delta_i/e_i
                     \leq (K_Y+D)\cdot A+\sum_i a_i.       \tag{D.UB.10}
\]
In particular, $s$ need not be isolated in $T$.
\end{proposition}

\begin{proof}
Repeat the resolution and contraction construction of the
preceding proposition. Resolution of the normal source may
also introduce curves over points outside the exceptional
locus, but these have zero intersection with $f^*A$.
For the point-mapped components in the support of $f^*A$
the same contractions and adjunction estimates apply.
Thus (D.UB.8) remains valid.

Decompose the relative canonical divisor into components
dominating exceptional curves, components mapped to points,
and the horizontal part $R_{\rm hor}$ dominating $D$.
Tameness at an original boundary divisor implies that at
every prolongation there its coefficient is $e-1$ and the
residue extension is separable. Hence the coefficient of
$f_*R_{\rm hor}$ at that component of $D$ is
\[
                    \sum_\lambda(e_\lambda-1)f_\lambda
                          =n-\sum_\lambda f_\lambda\leq n.
\]
There are no other horizontal ramification components.
Since $A$ is effective and exceptional, no component of $D$
is contained in $A$, and intersections are nonnegative.
Projection gives
\[
                 R_{\rm hor}\cdot f^*A
                      =(f_*R_{\rm hor})\cdot A\leq nD\cdot A.
\]
Add this term to the right side of (D.UB.9), then use
(D.UB.8). This gives (D.UB.10).
\end{proof}

\begin{corollary}[Proved: a rank-independent exceptional break bound]
\label{D:surface-uniform-conductor}
Assume that the ground field has positive characteristic.
Put $C_A=(K_Y+D)\cdot A+\sum_i a_i$, where $D=0$ in the
punctured case. Every finite cover in the preceding
propositions, and every finite-monodromy representation
arising from it, has nonlogarithmic and logarithmic
Abbes--Saito breaks at $E_i$ at most
\[
                         2\lceil C_A/c_i\rceil.           \tag{D.UB.11}
\]
For a sole smooth exceptional curve of genus $g\geq2$
and square $-d$ in the punctured case, the stronger
bound is
\[
                         2\left\lceil1+(2g-2)/d\right\rceil.
                                                               \tag{D.UB.12}
\]
No finite-generation conclusion is asserted from these
bounds alone.
\end{corollary}

\begin{proof}
Every $\delta_i$ is nonnegative. Apply the separate
normalized-different bounds and
Lemma~\ref{D:different-conductor-bound} to the completed
boundary DVRs. An arbitrary finite cover is dominated by
its Galois closure, which remains in the same cover class.
For logarithmic breaks use
$G^a_{\log}\subseteq G^a$ from
\cite[Proposition 3.15(1), preprint p.8]{AbbesSaito2002D}.
For the sole positive-genus curve, Galois conjugacy in
(D.PG.1) gives
\[
    \delta/e\leq1-1/e+
       \frac{2g-2}{d}\left(1-\frac1{f^{\rm insep}}\right)
                  \leq1+\frac{2g-2}{d}.
\]
A finite-monodromy representation factors through one
finite Galois cover, and none of these bounds involves
its dimension.
\end{proof}
